\documentclass[14pt,a4paper]{scrreprt}

\usepackage{cmap}
\usepackage[T1]{fontenc} 
\usepackage[utf8]{inputenc}
\usepackage[english,russian]{babel}

\usepackage{caption}
\usepackage{subcaption}

\usepackage{float}

\usepackage{enumitem}

\usepackage{graphicx}
\usepackage{multirow}


\usepackage{pgfplots}
\pgfplotsset{compat=newest}
\usepgfplotslibrary{units}

\usepackage{longtable}

\usepackage{caption}
\captionsetup{labelsep=endash}
\captionsetup[figure]{name={Рисунок}}
\captionsetup[subtable]{labelformat=simple}
\captionsetup[subfigure]{labelformat=simple}
\renewcommand{\thesubtable}{\text{Таблица }\arabic{chapter}\text{.}\arabic{table}\text{.}\arabic{subtable}\text{ --}}
\renewcommand{\thesubfigure}{\text{Рисунок }\arabic{chapter}\text{.}\arabic{figure}\text{.}\arabic{subfigure}\text{ --}}


\usepackage{textcomp}

\usepackage{amsmath}
\usepackage{amsfonts}
\usepackage{array}

\usepackage{geometry}
\geometry{left=30mm}
\geometry{right=15mm}
\geometry{top=20mm}
\geometry{bottom=20mm}
\geometry{foot=1.7cm}

\usepackage{titlesec}
\titleformat{\section}
{\normalsize\bfseries}
{\thesection}
{1em}{}
\titlespacing*{\chapter}{0pt}{-30pt}{8pt}
\titlespacing*{\section}{\parindent}{*4}{*4}
\titlespacing*{\subsection}{\parindent}{*4}{*4}

% Маркировка для списков
\def\labelitemi{$\circ$}
\def\labelitemii{$*$}


\usepackage{setspace}
\onehalfspacing % Полуторный интервал

\frenchspacing
\usepackage{indentfirst} % Красная строка

\usepackage{titlesec}
\usepackage{xcolor}
% Названия глав
\titleformat{\section}{\normalsize\textmd}{\thesection}{1em}{}

\definecolor{gray35}{gray}{0.35}

\newcommand{\hsp}{\hspace{20pt}} % длина линии в 20pt

\titleformat{\chapter}[hang]{\Huge}{\textcolor{gray35}{\thechapter.}\hsp}{0pt}{\Huge\textmd}

\titleformat{\section}{\Large}{\textcolor{gray35}\thesection}{20pt}{\Large\textmd}
\titleformat{\subsection}{\Large}{\thesubsection}{20pt}{\Large\textmd}
\titleformat{\subsubsection}{\normalfont\textmd}{}{0pt}{}

% Настройки введения

\addtocontents{toc}{\setcounter{tocdepth}{2}}
\addtocontents{toc}{\setcounter{secnumdepth}{1}}

\usepackage{tocloft,lipsum,pgffor}

\addtocontents{toc}{~\hfill\textnormal{Страница}\par}

\renewcommand{\cftpartfont}{\normalfont\textmd}

\addto\captionsrussian{\renewcommand{\contentsname}{Содержание}}
\renewcommand{\cfttoctitlefont}{\Huge\textmd}

\renewcommand{\cftchapfont}{\normalfont\normalsize}
\renewcommand{\cftsecfont}{\normalfont\normalsize}
\renewcommand{\cftsubsecfont}{\normalfont\normalsize}
\renewcommand{\cftsubsubsecfont}{\normalfont\normalsize}

\renewcommand{\cftchapleader}{\cftdotfill{\cftdotsep}}

\usepackage{listings}
\usepackage{xcolor}

\usepackage[pdftex]{hyperref} % Гиперссылки
\hypersetup{hidelinks}

% Листинги 
\usepackage{listings}

\definecolor{darkgray}{gray}{0.15}

\definecolor{teal}{rgb}{0.25,0.88,0.73}
\definecolor{gray}{rgb}{0.5,0.5,0.5}
\definecolor{b-red}{rgb}{0.88,0.25,0.41}
\definecolor{royal-blue}{rgb}{0.25,0.41,0.88}

\usepackage{listings}
\lstset{
	aboveskip=3mm,
	belowskip=3mm,
	frame=tb,
	frame=single,
	basicstyle=\footnotesize\ttfamily,
	numberstyle=\tiny\color{gray},
	keywordstyle=\color{royal-blue},
	commentstyle=\color{gray35},
	stringstyle=\color{b-red},
	numbers=left,
	numbersep=5pt,
	numberstyle=\tiny,
	showstringspaces=false, 
	captionpos=t,
	tabsize=4,
	language=Lisp
}

% какой то сложный кусок со стак эксчейндж для квадратных скобок
\makeatletter
\newenvironment{sqcases}{%
	\matrix@check\sqcases\env@sqcases
}{%
	\endarray\right.%
}
\def\env@sqcases{%
	\let\@ifnextchar\new@ifnextchar
	\left\lbrack
	\def\arraystretch{1.2}%
	\array{@{}l@{\quad}l@{}}%
}
\makeatother

% и для матриц
\makeatletter
\renewcommand*\env@matrix[1][\arraystretch]{%
	\edef\arraystretch{#1}%
	\hskip -\arraycolsep
	\let\@ifnextchar\new@ifnextchar
	\array{*\c@MaxMatrixCols c}}
\makeatother


\begin{document}

\begin{titlepage}
	
	\newgeometry{left=2cm, right=2cm, top=2.5cm, bottom=2.5cm}
	\fontsize{12pt}{12pt}\selectfont
	
	\noindent \begin{minipage}{0.13\textwidth}
		\includegraphics[width=\linewidth]{assets/bmstu-logo.png}
	\end{minipage}
	\noindent\begin{minipage}{0.85\textwidth}\centering
		\textbf{\textsc{Министерство науки и высшего образования Российской Федерации}}\\
		\textbf{\textsc{Федеральное государственное бюджетное образовательное 	учреждение высшего образования}}\\
		\textbf{\textsc{Московский государственный технический университет имени 	Н.Э.~Баумана}}\\
		\textbf{\textsc{(национальный исследовательский университет)}}\\
		\textbf{\textsc{(МГТУ им. Н.Э. Баумана)}}\\
	\end{minipage}
	
	\noindent\rule{18cm}{1.5pt}
	
	\vspace{8mm}
	
	\noindent\textnormal{ФАКУЛЬТЕТ}\hspace{5mm} \underline{\textnormal{~~~~~~~~~~~~~~~~~~«Информатика и системы управления»~~~~~~~~~~~~~~~~~~}} \newline\newline
	\textnormal{КАФЕДРА}\hspace{5mm} \underline{\textnormal{~~«Программное обеспечение ЭВМ и информационные технологии»~~}}
	\newline\newline
	\textnormal{НАПРАВЛЕНИЕ ПОДГОТОВКИ}\hspace{5mm} \underline{\textnormal{~~~~~~«09.03.04 Программная инженерия»~~~~~~~~}}
	
	\vspace{2.5cm}
	
	\begin{center}
		\Large\textbf{\textsc{ОТЧЕТ}}\\
		\Large\textbf{\textsc{ПО ЛАБОРАТОРНОЙ РАБОТЕ №18}}\\
	\end{center}
	
	\vspace{1cm}
	
	\noindent\textnormal{Название:} \hspace{15mm} \underline{\textnormal{~~~~~~Формирование и модификация списков на Prolog~~~~~~~}}\noindent
	
	\vspace{1.3cm}
	
	\noindent\textnormal{Дисциплина:} \hspace{10mm} \underline{\textnormal{~~~~~~Функциональное и логическое программирование~~~~~~}}\noindent
	
	\vspace{1.5cm}
	
	\noindent\textnormal{Студент} \hspace{17mm}
	\underline{\textnormal{{~~~~ИУ7-64Б~~~}}}
	\hspace{20mm}
	\underline{\textnormal{\hphantom{~~~~~~~~~~~~~~~~~~~~~~~~~~~}}} \hspace{14mm}
	\underline{\textnormal{~~~С. Д. Параскун~~~}}
	
	\vspace{2mm}
	\noindent\textnormal{\hphantom{Студент}} \hspace{23mm}\noindent
	\fontsize{8pt}{8pt}
	\textnormal{Группа}\hspace{40mm}\textnormal{Подпись, дата} \hspace{30mm}\noindent\textnormal{И. О. Фамилия}
	
	\vspace{0.5cm}
	
	\fontsize{12pt}{12pt}\selectfont
	\noindent\textnormal{Преподаватель} \hspace{52mm}
	\underline{\textnormal{\hphantom{~~~~~~~~~~~~~~~~~~~~~~~~~~~}}} \hspace{14mm}
	\noindent\underline{\textnormal{~~Н. Б. Толпинская~~}}
	
	\vspace{2mm}
	\noindent\textnormal{\hphantom{Студент}} \hspace{17mm}\noindent
	\fontsize{8pt}{8pt}
	\hphantom{Группа}\hspace{43mm}\textnormal{Подпись, дата} \hspace{30mm}\noindent\textnormal{И. О. Фамилия}
	
	\vspace{0.5cm}
	
	\fontsize{12pt}{12pt}\selectfont
	\noindent\textnormal{Преподаватель} \hspace{52mm}
	\underline{\textnormal{\hphantom{~~~~~~~~~~~~~~~~~~~~~~~~~~~}}} \hspace{14mm}
	\noindent\underline{\textnormal{~~Ю. В. Строганов~~}}
	
	\vspace{2mm}
	\noindent\textnormal{\hphantom{Студент}} \hspace{17mm}\noindent
	\fontsize{8pt}{8pt}
	\hphantom{Группа}\hspace{43mm}\textnormal{Подпись, дата} \hspace{30mm}\noindent\textnormal{И. О. Фамилия}
	
	\vspace{1cm}
	
	\fontsize{12pt}{12pt}\selectfont
	
	\begin{center}
		\vfill
		Москва, ~\the\year
		~г.
	\end{center}
	\restoregeometry

\end{titlepage}


\thispagestyle{empty}

\chapter{Практическое задание}

\section{Задание}

Используя хвостовую рекурсию, разработать эффективную программу, позволяющую:
\begin{enumerate}
	\item сформировать список из элементов числового списка, больших заданного значения;
	\item сформировать список из элементов, стоящих на нечетных позициях исходного списка (нумерация от 0);
	\item удалить заданный элемент из списка (один или все вхождения);
	\item преобразовать список в множество (можно использовать ранее разработанные процедуры).
\end{enumerate}
Убедиться в правильности результатов.

Для одного из вариантов вопроса и 1-го задания составить таблицу, отражающую конкретный порядок работы системы.

\section{Код программы}

\begin{lstlisting}[language=Prolog]
domains

	integerList = integer*.

predicates

	biggerList(integerList, integer, integerList).
	oddList(integerList, integerList).
	delElemFullInclude(integerList, integer, integerList).
	makeSet(integerList, integerList).

clauses

	biggerList([], _, []).
	biggerList([H|T], Num, [H|ResT]) :- H > Num, !, biggerList(T, Num, 
		ResT).
	biggerList([_|T], Num, Res) :- biggerList(T, Num, Res).
	
	oddList([], []).
	oddList([_, H|T], [H|ResT]) :- !, oddList(T, ResT).
	oddList([_|T], ResT) :- oddList(T, ResT).
	
	delElemFullInclude([], _, []).
	delElemFullInclude([H|T], Elem, [H|ResT]) :- H <> Elem, !, 
		delElemFullInclude(T, Elem, ResT).
	delElemFullInclude([_|T], Elem, Res) :- delElemFullInclude(T, Elem, 
		Res).
	
	makeSet([], []).
	makeSet([H|T], [H|ResT]) :- delElemFullInclude(T, H, Res), makeSet(Res, 
		ResT).

goal

	%biggerList([1, 7, 0, 4], 3, List).
	%oddList([1, 7, 0, 4], List).
	%delElemFullInclude([1, 6, 1, 0, 1], 1, List).
	makeSet([1, 5, 9, 1, 0, 5, 1], Set).
\end{lstlisting}

\section{Таблицы выполнения программы}

Запрос: biggerList([1, 7, 0, 4], 3, List).
\begin{table}[H]
	\begin{tabular}{|p{0.8cm\small}|p{4.7cm\small}|p{5.7cm\small}|p{4cm\small}|}	
		\hline
		№ шага & Состояния резольвенты и вывод: дальнейшие действия (почему?) & Для каких термов запускается алгоритм унификации: T1=T2 и каков результат (подстановка) & Дальнейшие действия: прямой ход или откат (почему и к чему приводит?)\\
		\hline
		1 & Резольвента: \begin{itemize}
			\item biggerList([1, 7, 0, 4], 3, List)
		\end{itemize} & biggerList([1, 7, 0, 4], 3, List) = biggerList([], \_, []). Унификация неуспешна & Прямой ход, переход к следующему предложению\\
		\hline
	\end{tabular}
\end{table}

\begin{table}[H]
	\begin{tabular}{|p{0.8cm\small}|p{4.7cm\small}|p{5.7cm\small}|p{4cm\small}|}	
		\hline
		№ шага & Состояния резольвенты и вывод: дальнейшие действия (почему?) & Для каких термов запускается алгоритм унификации: T1=T2 и каков результат (подстановка) & Дальнейшие действия: прямой ход или откат (почему и к чему приводит?)\\
		\hline
		2 & Резольвента: \begin{itemize} \item 1 > 3, \item !, \item biggerList([7, 0, 4], 3, ResT), \item ! \end{itemize} & biggerList([1, 7, 0, 4], 3, List) = biggerList([H|T], Num, [H|ResT]). Унификация успешна\linebreak \textbf{Подстановка:} \{H=1, T=[7, 0, 4], Num=3, List=[1|ResT]\} & Прямой ход, решение цели резольвенты\\
		\hline
		3 & Резольвента: \begin{itemize} \item !, \item biggerList([7, 0, 4], 3, ResT), \item ! \end{itemize} & 1 > 3. Ложь & Откат, переход к следующему относительно 2 предложению\\
		\hline
		4 & Резольвента: \begin{itemize} \item biggerList([7, 0, 4], 3, Res) \end{itemize} & biggerList([1, 7, 0, 4], 3, List) = biggerList([\_|T], Num, Res). Унификация успешна\linebreak \textbf{Подстановка:} \{T=[7, 0, 4], Num=3, List=Res\} & Прямой ход, решение цели резольвенты\\
		\hline
		5 & Резольвента: \begin{itemize} \item biggerList([7, 0, 4], 3, Res) \end{itemize} & biggerList([7, 0, 4], 3, Res) = biggerList([], \_, []). Унификация неуспешна & Прямой ход, переход к следующему предложению\\
		\hline
		6 & Резольвента: \begin{itemize} \item 7 > 3, \item !, \item biggerList([0, 4], 3, ResT), \item ! \end{itemize} & biggerList([7, 0, 4], 3, Res) = biggerList([H|T], Num, [H|ResT]). Унификация успешна\linebreak \textbf{Подстановка:} \{T=[7, 0, 4], Num=3, List=Res, H=7, T=[0, 4], Num=3, Res=[7|ResT]\} & Прямой ход, решение цели резольвенты\\
		\hline
		7 & Резольвента: \begin{itemize} \item !, \item biggerList([0, 4], 3, ResT), \item ! \end{itemize} & 7 > 3. Правда & Прямой ход, решение цели резольвенты\\
		\hline
		8 & Резольвента: \begin{itemize} \item biggerList([0, 4], 3, ResT), \item ! \end{itemize} & !. Отсечение 6, 7 & Прямой ход, решение цели резольвенты\\
		\hline
	\end{tabular}
\end{table}

\begin{table}[H]
	\begin{tabular}{|p{0.8cm\small}|p{4.7cm\small}|p{5.7cm\small}|p{4cm\small}|}	
		\hline
		№ шага & Состояния резольвенты и вывод: дальнейшие действия (почему?) & Для каких термов запускается алгоритм унификации: T1=T2 и каков результат (подстановка) & Дальнейшие действия: прямой ход или откат (почему и к чему приводит?)\\
		\hline
		9 & Резольвента: \begin{itemize} \item biggerList([0, 4], 3, ResT), \item ! \end{itemize} & biggerList([0, 4], 3, ResT) = biggerList([], \_, []). Унификация неуспешна & Прямой ход, переход к следующему предложению\\
		\hline
		10 & Резольвента: \begin{itemize}  \item 0 > 3, \item !, \item biggerList([4], 3, ResT), \item !, \item ! \end{itemize} & biggerList([0, 4], 3, ResT) = biggerList([H|T], Num, [H|ResT]). Унификация успешна\linebreak \textbf{Подстановка:} \{..., H=0, T=[4], Num=3, ResT=[0|ResT]\} & Прямой ход, решение цели резольвенты\\
		\hline
		11 & Резольвента: \begin{itemize}\item !, \item biggerList([4], 3, ResT), \item !, \item ! \end{itemize} & 0 > 3. Ложь & Откат, переход к следующему относительно 10 предложению\\
		\hline
		12 & Резольвента: \begin{itemize} \item biggerList([4], 3, Res), \item ! \end{itemize} & biggerList([0, 4], 3, ResT) = biggerList([\_|T], Num, Res). Унификация успешна\linebreak \textbf{Подстановка:} \{..., T=[4], Num=3, ResT=Res\} & Прямой ход, решение цели резольвенты\\
		\hline
		13 & Резольвента: \begin{itemize} \item biggerList([4], 3, Res), \item ! \end{itemize} & biggerList([4], 3, Res) = biggerList([], \_, []). Унификация неуспешна & Прямой ход, переход к следующему предложению\\
		\hline
		14 & Резольвента: \begin{itemize} \item 4 > 3, \item !, \item biggerList([], 3, ResT), \item !, \item ! \end{itemize} & biggerList([4], 3, Res) = biggerList([H|T], Num, [H|ResT]). Унификация успешна\linebreak \textbf{Подстановка:} \{..., H=4, T=[], Num=3, Res=[4|ResT]\} & Прямой ход, решение цели резольвенты\\
		\hline
	\end{tabular}
\end{table}

\begin{table}[H]
	\begin{tabular}{|p{0.8cm\small}|p{4.7cm\small}|p{5.7cm\small}|p{4cm\small}|}	
		\hline
		№ шага & Состояния резольвенты и вывод: дальнейшие действия (почему?) & Для каких термов запускается алгоритм унификации: T1=T2 и каков результат (подстановка) & Дальнейшие действия: прямой ход или откат (почему и к чему приводит?)\\
		\hline
		15 & Резольвента: \begin{itemize} \item !, \item biggerList([], 3, ResT), \item !, \item ! \end{itemize} & 4 > 3. Правда & Прямой ход, решение цели резольвенты\\
		\hline
		16 & Резольвента: \begin{itemize} \item biggerList([], 3, ResT), \item !, \item ! \end{itemize} & !. Отсечение 14, 15 & Прямой ход, решение цели резольвенты\\
		\hline
		17 & Резольвента: \begin{itemize} \item !, \item ! \end{itemize} & biggerList([], 3, ResT) = biggerList([], \_, []). Унификация успешна\linebreak \textbf{Подстановка:} \{..., H=4, T=[], Num=3, Res=[4|[]], ResT=[]\} & Прямой ход, решение цели резольвенты\\
		\hline
		18 & Резольвента: \begin{itemize} \item ! \end{itemize} & !. Отсечение 14-16 & Откат, переход к шагу 14\\
		\hline
		19 & Резольвента: пуста\linebreak \textbf{Вывод:} List=[7, 4] & !. Отсечение 6-13 & Откат, переход к шагу 6, завершение работы\\
		\hline
	\end{tabular}
\end{table}
		
\end{document}